\section{无穷级数与部分和}
\label{sec:02-Calculus/05-Sequences-and-Series/02}

\subsection{无穷多项相加，不能一个一个来}

无穷级数的符号
\[
\sum_{n=1}^{\infty}a_n
\]
看上去像是在说：把 \(a_1,a_2,a_3,\dots\) 全部加起来。但如果你真的试图一项一项往下加，你永远做不完——加法是一个二元运算，它在任意有限步内只能处理有限多个数。无穷个数的``总和''这个概念本身就需要被重新发明。

这就是无穷级数与普通加法的根本张力：你脑子里想做的是``无限次加法''，但数学里能操作的只有有限。需要一个桥把两者连起来。

\subsection{部分和：有限的桥，无穷的极限}

出路是把无穷和转化为我们已经理解的东西——数列极限。

先截断到第 \(N\) 项，算一个有限的累加：
\[
S_N=\sum_{n=1}^N a_n,
\]
称为第 \(N\) 部分和。\(S_1=a_1\)，\(S_2=a_1+a_2\)，\(S_3=a_1+a_2+a_3\)……随着 \(N\) 增大，\(S_N\) 逐渐靠近我们心中那个``总和''的位置，如果它确实在靠近某个数的话。

于是无穷级数的收敛就定义为：部分和数列 \((S_N)\) 收敛。若
\[
\lim_{N\to\infty}S_N=S,
\]
则称级数收敛，和为 \(S\)；否则发散。

这个定义的精妙之处在于，它从不要求你做无限次加法。每一步都是有限的 \(S_N\)，最后只对一个数列取一次极限。无穷和不是加法运算的延伸——它是有限和序列的极限行为。前面有限项改什么都不会影响尾部的收敛性，它只关心 \(N\) 越来越大时 \(S_N\) 的行为。

\subsection{几何级数：不用猜，部分和封闭可算}

几何级数
\[
1+r+r^2+r^3+\cdots
\]
是少数能直接写出部分和封闭公式的级数之一。考虑
\[
S_N=1+r+r^2+\cdots+r^N.
\]

把 \(S_N\) 乘以 \(r\)：
\[
rS_N=r+r^2+r^3+\cdots+r^{N+1}.
\]

两式相减：
\[
S_N-rS_N=1-r^{N+1},
\]
所以当 \(r\neq1\) 时
\[
S_N=\frac{1-r^{N+1}}{1-r}.
\]

现在取极限。若 \(|r|<1\)，\(r^{N+1}\to0\)，于是
\[
\sum_{n=0}^{\infty}r^n=\lim_{N\to\infty}S_N=\frac1{1-r}.
\]

若 \(|r|\ge1\) 且 \(r\neq1\)，\(r^{N+1}\) 不趋于零（绝对值不小于 1），极限不存在。若 \(r=1\)，原级数是 \(1+1+1+\cdots\)，部分和 \(S_N=N+1\to\infty\)，发散。若 \(r=-1\)，级数是 \(1-1+1-1+\cdots\)，部分和交替为 1 和 0，不收敛。

一个可以挂在嘴边的例子：无限小数 \(0.999\ldots\) 即
\[
\frac9{10}+\frac9{100}+\frac9{1000}+\cdots
= \frac{9/10}{1-1/10}=1.
\]
它不是一个``永远比 1 差一点点''的未完成过程，而是几何级数极限恰好等于 1。

\subsection{项趋于零：门槛最低的必要条件}

若 \(\sum a_n\) 收敛，令 \(S_n\) 为前 \(n\) 项的和，则
\[
a_n=S_n-S_{n-1}\to S-S=0.
\]

这就是说：级数收敛的必要条件是通项趋于零。逻辑上相当于：如果连这个条件都不满足，级数已经判了死刑——不需要进一步分析。比如 \(\sum n/(n+1)\)，通项趋于 1，立刻断定发散。

但反向不成立。通项趋于零只是门票，不保证能进场。

\subsection{调和级数：教科书级反例}

\[
\sum_{n=1}^{\infty}\frac1n
=1+\frac12+\frac13+\frac14+\frac15+\cdots
\]

通项 \(1/n\to0\)，所以必要条件通过了。但部分和却发散。

我们来亲手验证它为什么发散。把项按长度翻倍的方式分组：
\[
1+\frac12+
\underbrace{\left(\frac13+\frac14\right)}_{\textstyle 2\text{项}}+
\underbrace{\left(\frac15+\frac16+\frac17+\frac18\right)}_{\textstyle 4\text{项}}+
\underbrace{\left(\frac19+\cdots+\frac1{16}\right)}_{\textstyle 8\text{项}}+\cdots
\]

看第三组（\(k\ge3\)，块大小 \(2^{k-1}\) 项）：块内每一项都大于或等于块的末项。对于第 \(k\) 组，末项是 \(1/2^k\)，块的大小是 \(2^{k-1}\)，所以块的总和至少是
\[
2^{k-1}\cdot\frac1{2^k}=\frac12.
\]

从第三组开始，每一大块至少贡献 \(1/2\)。有无穷多个这样的块，所以部分和可以超过任何有限上界。注意发生了什么：项确实越来越小，但项数增长的速度抵消了衰减的速度。调和级数不是在某一项变得不够小——它的问题在于``变小得不够快''。

这种局面比直觉更难察觉。考虑一个长期运行的服务器：每次处理第 \(n\) 个请求时，程序泄漏 \(1/n\) MB 内存。从监控图看，最新一次泄漏已经小到可以忽略。但处理完 \(N\) 个请求后的总泄漏是 \(\sum_{n=1}^N 1/n \sim \ln N\)——没有上界。``单次增量在减小''和``累计状态已经受控''是两个完全不同的命题，后者需要检查的是部分和，不是通项。

\subsection{望远镜级数：中间项全部消去的幸运结构}

有些级数的部分和可以写成相消的形式：
\[
a_n=\frac1n-\frac1{n+1},
\]
则
\[
S_N=\left(1-\frac12\right)+\left(\frac12-\frac13\right)+\cdots+\left(\frac1N-\frac1{N+1}\right)
=1-\frac1{N+1}\to1.
\]

所有中间项成对消失，只剩首尾。但这里有一个容易踩的陷阱：你不能在无穷级数的符号里直接``首尾相消''就宣称答案——必须先把第 \(N\) 部分和的消去结果写清楚，看到剩下的边界项在取极限后确实趋于零。否则可能漏掉一个不会消失的残余项。

比如若写成
\[
\sum_{n=1}^{\infty}\left(\frac1n-\frac1{n+2}\right),
\]
前几项展开看：\((1-1/3)+(1/2-1/4)+(1/3-1/5)+(1/4-1/6)+\cdots\)，消去后留下 \(1+1/2=3/2\) 而不是简单的首项。部分和的正确展开会发现有两个边界项存活。

\subsection{Cauchy 准则：不依赖极限值的收敛判断}

收敛的定义要求你预先知道或猜出极限 \(S\)，然后验证 \(S_N\) 是否趋近于它。但如果极限值不好猜呢？Cauchy 准则提供了一个只看序列内部行为的方法：

\[
\sum a_n\text{ 收敛}
\Longleftrightarrow
\forall\varepsilon>0,\;\exists N,\;\forall m>n\ge N,\;
|a_{n+1}+a_{n+2}+\cdots+a_m|<\varepsilon.
\]

翻译过来：对任意给定的容差，存在一个位置 \(N\)，使得从这个位置之后截取的任意有限长度的尾段，累加和都小于容差。也就是说，足够远的尾部，无论怎么截取和累加，贡献都微不足道。

这个准则不需要你预先知道级数的和。它直接抓住的是``尾部足够轻''这个直觉，并且把它变成一个可操作的不等式条件。从数列的角度看，Cauchy 准则就是说部分和序列 \((S_N)\) 是一个 Cauchy 列——而实数系的完备性保证每个 Cauchy 列都收敛。所以级数的收敛问题，追到底，是实数完备性在起作用。

Cauchy 准则是绝对收敛、一致收敛等后续概念的共同原型——它关心的就是尾部能不能被任意地压扁。

\subsection{记住这个不对称性}

级数收敛 \(\\Longrightarrow\) 项趋于零，但反过来不成立。调和级数把不对称性展示得清清楚楚。

把几何级数 \(1+1/2+1/4+\cdots\) 和调和级数 \(1+1/2+1/3+\cdots\) 并排放在一起对比最能说明问题。前者的项按公比 \(1/2\) 缩小，每前进一步项就减半；后者的项也趋于零，但缩水速度是 \(1/(n+1)\approx 1-1/n\)，几乎不变。几何衰减和代数衰减之间的差距，就是收敛和发散之间的那道门槛。

在后面几节里你会看到，判断级数是否收敛的任务，很大程度上就是在项已经趋于零的前提下，进一步确认尾部的累积是否真的被控在一个有限量之内。对正项级数这件事变得相对可操作——因为部分和只能朝一个方向走，下一节就从这个有利地形出发。
